% SGA 6, Exposición X: fragmento español alineado con la fuente.
% Alcance: Ejemplo 1.5 completo, inmediatamente antes de la subsección 1.6.
% Fuente francesa corregida: sga6/split_r2/corr/main.tex, líneas 12631--12670.
% SHA-256 de la fuente completa: 5D841302BF42D65817A80CC3754987AE3798A3E0D0DA2207177BE9A4C21B24AF.
% Control diplomático: el bloque correspondiente de orig/main.tex es idéntico.
% Integración: requiere tikz-cd y las macros \codim, \dimdot y \degtr.

\textbf{Ejemplo 1.5.} Sea $S$ un esquema localmente noetheriano
universalmente catenario (por ejemplo, un subesquema de un esquema de
Cohen--Macaulay). Definamos
\begin{equation}\tag{1.5.1}
D_S(x)=n-\codim(\overline{\{x\}},S)=n-\dimdot\mathcal O_{S,x}\qquad(x\in S),
\end{equation}
y, para todo $S$-esquema propio $s_Y:Y\to S$ y todo $y\in Y$, definamos
también
\begin{equation}\tag{1.5.2}
D_Y(y)=D_S(s_Y(y))+\degtr_{k(s_Y(y))}k(y).
\end{equation}
Si $f:Y_1\to Y_2$ es un $S$-morfismo de $S$-esquemas propios, la aditividad
de los grados de trascendencia y (1.5.2) dan inmediatamente
\begin{equation}\tag{1.5.3}
D_{Y_1}(y)=D_{Y_2}(f(y))+\degtr_{k(f(y))}k(y),\qquad(y\in Y_1).
\end{equation}
Mostremos que las funciones $D_Y$ así definidas satisfacen las condiciones
(i)--(iv) de 1.4. En virtud de (1.5.3), la condición (ii) es la única no
trivial. Gracias a (iv), podemos suponer evidentemente que
$Y=\overline{\{x\}}$. Factorizamos entonces el morfismo estructural $s_Y$ de
$Y$ a través de su imagen $s_Y(Y)=Z\hookrightarrow S$; es decir,
$s_Y=g\circ f$:
\[
\begin{tikzcd}
Y \arrow[r,"f"] & Z \arrow[r,hook,"g"] & S,
\end{tikzcd}
\]
donde $Y$ y $Z$ son irreducibles, $f$ es propio y sobreyectivo, y $g$ es una
inmersión cerrada. Pongamos $x'=s_Y(x)$, $y'=s_Y(y)$. Se deduce entonces de
(1.5.2) y (1.5.1) que
\[
D_Y(x)-D_Y(y)=\dimdot\mathcal O_{S,y'}-\dimdot\mathcal O_{S,x'}+
\degtr_{k(x')}k(x)-\degtr_{k(y')}k(y).
\]
Puesto que $S$ es catenario, se tiene
\[
\dimdot\mathcal O_{S,y'}-\dimdot\mathcal O_{S,x'}=
\dimdot\mathcal O_{Z,y'}-\dimdot\mathcal O_{Z,x'}
\]
(EGA $0_{IV}$ 14.3.2.1), de donde
\begin{equation}\tag{1.5.4}
D_Y(x)-D_Y(y)=
[\dimdot\mathcal O_{Z,y'}-\degtr_{k(y')}k(y)]-
[\dimdot\mathcal O_{Z,x'}-\degtr_{k(x')}k(x)].
\end{equation}
Como $Z$ es universalmente catenario, se puede aplicar EGA IV 5.6.5.1, y se
obtiene
\[
C_1=\dimdot\mathcal O_{Y,y}-e,\qquad C_2=\dimdot\mathcal O_{Y,x}-e,
\]
donde $C_1$ y $C_2$ son los dos términos del miembro derecho de (1.5.4).
Se concluye entonces inmediatamente de (1.5.4) que
\[
D_Y(x)-D_Y(y)=\dimdot\mathcal O_{Y,y}-\dimdot\mathcal O_{Y,x}>0,
\qquad \text{C.Q.D.}
\]
